\documentclass[10pt,a4paper]{amsart}
\usepackage[margin=19mm]{geometry}
\usepackage{amsmath,amssymb,mathtools}
\usepackage[scr=rsfs,cal=euler]{mathalfa}
\usepackage{xfrac}
\usepackage[unicode,hidelinks,bookmarks=false]{hyperref}
\newcommand{\ie}{i.e.\ }
\newcommand{\cf}{cf.\ }
\newcommand{\resp}{respectively\ }
\newcommand{\Subsection}{\subsection}
\providecommand{\nobreakdash}{\nobreak\hbox{-}}
\setlength{\emergencystretch}{2em}
\multlinegap0pt
\usepackage{amssymb}
\usepackage[shortlabels]{enumitem}
\usepackage{xcolor}
\newtheorem{lemma}{Lemma}
\newtheorem{remark}{Remark}
\newtheorem{theorem}{Theorem}
\newcommand{\R}{\mathbb{R}}
\newcommand{\X}{\mathcal{X}}
\newcommand{\Y}{\mathcal{Y}}
\newcommand{\qbox}[1]{\quad\hbox{#1}\quad}
\title{Inertial Primal Dual Dynamics with Hessian-driven Damping for Saddle Point Problems}
\author{Zepeng Wang, Juan Peypouquet}
\date{Source: arXiv:2607.26235v2 (CC BY 4.0)}
\begin{document}
\maketitle

\noindent\emph{Source and licence.} Inertial Primal Dual Dynamics with Hessian-driven Damping for Saddle Point Problems. Version arXiv:2607.26235v2 (CC BY 4.0), distributed by arXiv under the Creative Commons Attribution 4.0 International licence, http://creativecommons.org/licenses/by/4.0/, as recorded in the arXiv licence metadata for this exact version and shown on the abstract page https://arxiv.org/abs/2607.26235v2. Full licence notice: \texttt{licenses/arxiv-2607.26235-CC-BY-4.0.txt}.\par\smallskip

\noindent\textit{Editorial scope.} Counted unit: the article's theorem T:rates (source0.tex lines 982--991). The article's complete original proof of it is delivered with it (source0.tex lines 994--1069, one \texttt{proof} environment of 76 lines). No mathematics is written, repaired or re-derived: the statement and the proof are the article's own lines, in the article's own notation, and no retained line is substituted or re-flowed.  Retained uncounted as the source-local context the proof needs: lines 755--763; lines 800--806; lines 809--819; lines 899--904; lines 935--943; lines 948--951; lines 958--965.  Nothing was omitted from any retained span, so the package carries no omission record.  No retained line contains a citation, so the wrapper carries no bibliography and no bibliography entry.  No retained line contains a figure environment, an image-inclusion command or drawing code, so no image is delivered and the package declares no asset dependency.  Editorial formatting only, in addition: an AMS \texttt{amsart} preamble that restates the article's own declarations and macros used by the retained lines, namely the environments \texttt{lemma}, \texttt{remark}, \texttt{theorem} on separate counters, together with the standard packages those lines need.  The retained environments print in this document as Lemma 1, Remark 1, Theorem 1 in order of appearance, whereas the article prints the same items with the numbering of its own full document.  The complete native source of the article as distributed in the arXiv e-print for this exact version is bundled unchanged as \texttt{sources/206248/source0.tex}.
\begin{equation}\label{inertial_primal_dual_system}\tag{PD-AVD-H}
\left\{
\begin{array}{rcl}
\ddot{x} + \frac{\alpha}{t}\dot{x} + \beta\nabla^2_{xx}\mathcal{L}_\rho \left( x, y \right) \dot{x} 
 + \left( \gamma + \frac{r}{t} \right) \nabla_x \mathcal{L}_\rho\left( x, y + \theta t \dot{y} \right) &=& 0, \\[2pt]
\ddot{y} + \frac{\alpha}{t}\dot{y} - \left( \gamma + \frac{r}{t} \right) \nabla_y \mathcal{L}_\rho\left( x + \theta t \dot{x}, y \right) &=& 0,
\end{array}
\right.
\end{equation}

\begin{equation}\label{E: W}
\begin{aligned}
W &= \frac{1}{2}\| t\dot{s} + \xi(s-s^*) \|^2 + \frac{\eta}{2}\| s-s^* \|^2 + t(\gamma t + r) F(x) \\ 
&\quad + \beta\xi t \left[ \langle \nabla f(x), x-x^* \rangle - \left( f(x) - f(x^*) \right) \right] 
 + \frac{\beta\xi \rho t}{2}\| Ax-b \|^2,
\end{aligned}
\end{equation}

\begin{lemma}\label{Lem: epsilon_dev}
Let $W$ be defined by \eqref{E: W}. Set $\xi = \frac{1}{\theta}$. For every $t\ge t_0$, we have
\begin{align*}
\dot{W}
&= -(\alpha-\xi-1)t\| \dot{s} \|^2
   -\beta t^2 \left\langle \left( \nabla^2 f(x) + \rho A^* A \right)\dot{x}, \dot{x} \right\rangle \\
&\quad - \frac{\rho\xi}{2}( \gamma t + r - \beta ) \| Ax-b \|^2   
  - [ (\xi-2)\gamma t + (\xi-1)r ] F(x) \\
&\quad -\xi( \gamma t + r - \beta ) \left[ \langle \nabla f(x), x-x^* \rangle - \left( f(x) - f(x^*) \right) \right].
\end{align*}
\end{lemma}

\begin{remark}\label{Rem: epsilon_dev_bound}
Since $f$ is convex, we have 
$$ \left\langle \big( \nabla^2 f(x) + \rho A^* A \big)\dot{x}, \dot{x} \right\rangle 
= \left\langle \nabla^2 f(x) \dot{x}, \dot{x} \right\rangle + \rho \| A\dot{x} \|^2 \ge 0.$$
If $\alpha\ge 3$, $r\ge\beta$ and $\frac{1}{\alpha-1}\le \theta \le \frac{1}{2}$, then $W(t)$ is nonincreasing, whence $W(t)\le W(t_0)$ for all $t\ge t_0$.
\end{remark}

\begin{equation}\label{E: epsilon_bar}
\begin{aligned}
E_\ell
&= \frac{1}{2}\| t\dot{s}(t) + \xi(s(t)-s_\ell^*) \|^2 + \frac{\eta}{2}\| s(t)-s_\ell^* \|^2
  + t(\gamma t + r) \big[ \mathcal{L}_\rho\left(x(t),\ell\right) - \mathcal{L}_\rho(x^*,\ell)\big] \\
&\quad + \beta\xi t \left[ \langle \nabla f(x(t)), x(t)-x^* \rangle - \left( f(x(t)) - f(x^*) \right) \right] 
  + \frac{\beta\xi \rho t}{2}\| Ax(t)-b \|^2,
\end{aligned}
\end{equation}

\begin{equation}\label{E: Lagrangian_bar}
\mathcal{L}_\rho\left(x,\ell\right) - \mathcal{L}_\rho(x^*,\ell) = F(x) 
  + \left\langle \ell - y^*, Ax-b \right\rangle,
\end{equation}

\begin{lemma}\label{Lem: epsilon_bar}
Let $(x,y):[t_0,\infty)\to\X\times\Y$ be a solution of \eqref{inertial_primal_dual_system}. For $\ell\in\Y$, let $E_\ell:[t_0,\infty)\to\R$ be defined by \eqref{E: epsilon_bar}. Then, we have
\begin{equation*}
E_\ell \le 2W + (\xi^2 + \eta)\| \ell - y^*\|^2 + t(\gamma t + r)\left\langle \ell - y^*, Ax(t)-b \right\rangle,
\end{equation*}
for every $t\ge t_0$. In particular,
$$\sup_{\|\ell-y^*\|\le 1} E_\ell(t_0) \le C_0:=2W(t_0) + \xi^2 + \eta + t_0(\gamma t_0 + r) \| Ax_0-b \|.$$
\end{lemma}

\begin{theorem} \label{T:rates}
Let $(x,y):[t_0,\infty)\to\X\times\Y$ be a solution of the system \eqref{inertial_primal_dual_system}, where
$$ \alpha\ge 3,\quad r\ge\beta \qbox{and} \frac{1}{\alpha-1}\le \theta \le \frac{1}{2}.$$
Let $C_0$ be the constant defined in Lemma \ref{Lem: epsilon_bar}. For every $t\ge t_0$, we have
\begin{equation} \label{E:main_thm}
    0\le \mathcal{L}_\rho\left(x(t),y^*\right) - \mathcal{L}_\rho(x^*,y^*) + \|Ax(t)-b\| \le \frac{2C_0}{t(\gamma t + r)},
\end{equation}
and
$$ -\frac{2C_0 \| y^* \| }{t(\gamma t + r)} \le f\left( x(t) \right) - f(x^*) \le \frac{2C_0 \left(1+\|y^* \|\right)}{t(\gamma t + r)}. $$
\end{theorem}

\begin{proof}
Fix $\tau\ge t_0$, and set
\begin{equation}\label{E: lambda_tau}
\ell_\tau = \left\{
\begin{array}{ccl}
y^* + \frac{Ax(\tau)-b}{\| Ax(\tau)-b \|},& & \text{if }Ax(\tau)-b\neq 0,\\
y^* & & \text{if }Ax(\tau)-b= 0.
\end{array}
\right.
\end{equation}
Since $\|\ell_\tau-y^*\|\le 1$, it follows from Lemma \ref{Lem: epsilon_bar} that 
$$ E_{\ell_\tau}( t_0) \le C_0,$$
and
\begin{equation}\label{E: E_tau_bound}
\begin{aligned}
E_{\ell_\tau}(t) & \le 2W(t) + \xi^2 + \eta + t(\gamma t + r)\left\langle \ell_\tau - y^*, Ax(t)-b \right\rangle \\
 & \le 2W(t_0) + \xi^2 + \eta + t(\gamma t + r)\left\langle \ell_\tau - y^*, Ax(t)-b \right\rangle,    
\end{aligned}
\end{equation}
in view of Remark \ref{Rem: epsilon_dev_bound}. On the other hand, setting $\xi = \frac{1}{\theta}$, and proceeding as in the proof of Lemma \ref{Lem: epsilon_dev}, we obtain
\begin{align*}
\dot E_{\ell_\tau}(t)  
&\le -[ (\xi-2)\gamma t + (\xi-1)r ] \left[ \mathcal{L}_\rho\left(x(t),\ell_\tau\right) - \mathcal{L}_\rho(x^*,\ell_\tau) \right] \\
&\le -[ (\xi-2)\gamma t + (\xi-1)r ] \left\langle \ell_\tau - y^*, Ax(t)-b \right\rangle,
\end{align*}
where the last inequality is due to \eqref{E: Lagrangian_bar}. Multiplying both sides of this inequality by $\frac{t^{\xi-1}}{\gamma t+r}$ gives
\begin{align*}
\left(\tfrac{t^{\xi-1}}{\gamma t+r}\right)\dot E_{\ell_\tau}(t) 
& \le -\tfrac{t^{\xi-1}}{\gamma t+r}[ (\xi-2)\gamma t + (\xi-1)r ] \left\langle \ell_\tau - y^*, Ax(t)-b \right\rangle \\
& = -t(\gamma t+r) \left\langle \ell_\tau - y^*, Ax(t)-b \right\rangle \frac{d}{dt}\left(\tfrac{t^{\xi-1}}{\gamma t+r}\right).
\end{align*}
It follows that
\begin{align*}
\frac{d}{dt}\left[\left(\tfrac{t^{\xi-1}}{\gamma t+r}\right)E_{\ell_\tau}(t)\right] 
&= E_{\ell_\tau}(t)\frac{d}{dt}\left(\tfrac{t^{\xi-1}}{\gamma t+r}\right) + \left(\tfrac{t^{\xi-1}}{\gamma t+r}\right)\dot E_{\ell_\tau}(t)  \\
&\le \left[E_{\ell_\tau}(t)-t(\gamma t+r) \left\langle \ell_\tau - y^*, Ax(t)-b \right\rangle\right]\frac{d}{dt}\left(\tfrac{t^{\xi-1}}{\gamma t+r}\right) \\
&\le \left[2W(t_0) + \xi^2 + \eta \right] \frac{d}{dt}\left(\tfrac{t^{\xi-1}}{\gamma t+r}\right),
\end{align*}
by applying \eqref{E: E_tau_bound}. Integrating from $t_0$ to $t$, we obtain
\begin{align*}
\left(\tfrac{t^{\xi-1}}{\gamma t+r}\right)E_{\ell_\tau}(t)-\left(\tfrac{t_0^{\xi-1}}{\gamma t_0+r}\right)E_{\ell_\tau}(t_0) 
&\le  \left[2W(t_0) + \xi^2 + \eta \right] \left(\tfrac{t^{\xi-1}}{\gamma t+r}-\tfrac{t_0^{\xi-1}}{\gamma t_0+r}\right),
\end{align*}
which implies that
$$E_{\ell_\tau}(t) \le E_{\ell_\tau}(t_0) + \left[2W(t_0) + \xi^2 + \eta \right] \le 2C_0.$$
According to the definition of $E_{\ell}$ in \eqref{E: epsilon_bar}, we have
$$ \mathcal{L}_\rho\left(x(t),\ell_\tau\right) - \mathcal{L}_\rho(x^*,\ell_\tau) \le \frac{2C_0}{t(\gamma t + r)},$$
for every $t\ge t_0$. This, combined with \eqref{E: Lagrangian_bar} and \eqref{E: lambda_tau}, results in
$$ \mathcal{L}_\rho\left(x(\tau),y^*\right) - \mathcal{L}_\rho(x^*,y^*) + \| Ax(\tau)-b \|
\le \frac{2C_0}{\tau(\gamma \tau + r)}. $$
Since $\tau\ge t_0$ was arbitrarily chosen, we conclude that
$$ \mathcal{L}_\rho\left(x(t),y^*\right) - \mathcal{L}_\rho(x^*,y^*) + \| Ax(t)-b \| 
\le \frac{2C_0}{t(\gamma t + r)}, $$
for every $t\ge t_0$, which is \eqref{E:main_thm}. Now, \eqref{E:main_thm} also shows that
$$ \| Ax(t)-b \| \le \frac{2C_0}{t(\gamma t + r)}$$
and
$$f\left( x(t) \right) - f(x^*) + \left\langle y^*, Ax(t) - b \right\rangle 
\le \frac{2C_0}{t(\gamma t + r)}.$$
Hence,
$$ f\left( x(t) \right) - f(x^*)  
\le \frac{2C_0}{t(\gamma t + r)} + \|y^* \| \left\| Ax(t) - b \right\| 
\le \frac{2C_0 \left(1+\|y^* \|\right)}{t(\gamma t + r)}. $$
On the other hand, by the convexity of $f$ and the optimality conditions, we have
\begin{align*}
f\left( x(t) \right) - f(x^*) 
&\ge \langle \nabla f(x^*), x(t) - x^* \rangle 
= - \langle A^* y^*, x(t) - x^* \rangle 
= - \langle y^*, Ax(t) - b \rangle.
\end{align*}
This results in
$$ f\left( x(t) \right) - f(x^*) 
\ge - \| y^* \|  \left\| Ax(t) - b \right\| 
\ge -\frac{2C_0 \| y^* \| }{t(\gamma t + r)},
$$
and completes the proof.
\end{proof}

\end{document}
